\begin{tikzpicture}[>=Latex,scale=0.92,every node/.style={font=\scriptsize}]
  % Source and collimator
  \draw[rounded corners=2pt,fill=gray!12,thick] (-4.4,-0.55) rectangle (-3.45,0.55);
  \node at (-3.92,0.18) {$\alpha$};
  \node at (-3.92,-0.18) {源};
  \draw[fill=black] (-3.25,-0.42) rectangle (-3.10,0.42);
  \draw[fill=white] (-3.26,-0.07) rectangle (-3.09,0.07);
  \node[below] at (-3.18,-0.55) {准直孔};

  % Gold foil
  \draw[fill=yellow!25,draw=orange!80!black,thick] (-0.22,-1.40) rectangle (0.00,1.40);
  \node[above] at (-0.10,1.45) {金箔};

  % Fluorescent screen arc
  \draw[green!45!black,thick] (2.55,-1.75) arc[start angle=-37,end angle=37,radius=3.20];
  \node[green!35!black,align=center] at (3.80,0) {荧光屏\\记录闪光};

  % Incoming alpha particles
  \foreach \y in {-0.70,-0.35,0,0.35,0.70} {
    \draw[blue,thick,->] (-3.05,\y) -- (-0.24,\y);
  }
  \node[blue] at (-1.75,1.00) {高速 $\alpha$ 粒子};

  % Main paths
  \draw[blue,very thick,->] (0.00,0.70) -- (2.80,0.72);
  \draw[blue,very thick,->] (0.00,0.35) .. controls (1.00,0.46) and (1.80,0.85) .. (2.58,1.32);
  \draw[blue,very thick,->] (0.00,-0.35) .. controls (0.85,-0.48) and (1.55,-0.95) .. (2.20,-1.45);
  \draw[red,very thick,->] (0.00,0.02) .. controls (-0.55,0.46) and (-0.90,0.95) .. (-1.38,1.20);

  \node[align=center] at (1.38,-2.05) {绝大多数几乎直穿\\少数偏转，极少数大角度反弹};

  % Zoomed atom/nucleus
  \begin{scope}[shift={(0.20,2.15)},scale=0.78]
    \draw[rounded corners=2pt,gray] (-1.10,-0.82) rectangle (2.35,0.92);
    \node[above] at (0.62,0.92) {核式结构解释};
    \draw[gray,dashed] (0,0) circle (0.62);
    \fill[red] (0,0) circle (2.5pt);
    \node[red,right] at (0.10,0) {带正电的小原子核};
    \draw[blue,thick,->] (-0.95,-0.48) .. controls (-0.22,-0.35) and (-0.08,-0.12) .. (-0.88,0.48);
    \node[blue,left] at (-0.88,0.48) {$\alpha$};
  \end{scope}

  % Conclusions
  \node[align=left] at (-2.00,-1.70)
    {实验事实：\\原子大部分空间空旷；\\正电荷和质量集中在很小的核内。};
\end{tikzpicture}
